\mxnewpage\section*{EV1-B -- DS commun}\label{EV1-B}
\addcontentsline{toc}{section}{EV1-B -- DS commun}
\textbf{Sujet réservé au professeur}. Chapitres 1 à 4.\par Durée : 2 h. Sans calculatrice. Note sur 20. 14 novembre 2026.
\subsection*{Partie 1 -- Automatismes (6 points)}
Douze questions, \num{0,5} point chacune. Une seule réponse correcte ; aucune justification demandée. Une réponse incorrecte ou absente rapporte 0, sans retrait de points.
\begin{enumerate}
\item Calculer $6+2\times3$.\par A.~$11$\quad B.~$12$\quad C.~$15$\quad D.~$24$
\item Calculer $\frac12+\frac13$.\par A.~$\frac{1}{6}$\quad B.~$\frac{5}{6}$\quad C.~$\frac{2}{5}$\quad D.~$\frac{5}{3}$
\item Calculer $\dfrac{2^6}{2^{4}}$.\par A.~$2$\quad B.~$11$\quad C.~$1$\quad D.~$4$
\item Une valeur de 50 augmente de \pc{20}. Quelle est la valeur finale ?\par A.~$60$\quad B.~$70$\quad C.~$10$\quad D.~$40$
\item Une valeur de 80 est obtenue après une baisse de \pc{20}. Quelle est la valeur initiale ?\par A.~$64$\quad B.~$96$\quad C.~$60$\quad D.~$100$
\item Quel est le coefficient multiplicateur global d'une hausse de \pc{25} suivie d'une baisse de \pc{20} ?\par A.~$\frac{21}{20}$\quad B.~$\frac{4}{5}$\quad C.~$1$\quad D.~$\frac{9}{10}$
\item La solution de l'équation $2x+6=0$ est~:\par A.~$3$\quad B.~$12$\quad C.~$-3$\quad D.~$-12$
\item La pente de la droite passant par $\pt{0;6}$ et $\pt{2;12}$ est~:\par A.~$-3$\quad B.~$3$\quad C.~$6$\quad D.~$2$
\item La moyenne de $6, 8, 10$ est~:\par A.~$8$\quad B.~$24$\quad C.~$6$\quad D.~$9$
\item Un secteur de $90^\circ$ représente quelle fraction du total ?\par A.~$\frac{1}{3}$\quad B.~$\frac{1}{2}$\quad C.~$\frac{1}{4}$\quad D.~$\frac{3}{4}$
\item Une probabilité vaut $\dfrac{2}{5}$. Quelle est la probabilité de son complément ?\par A.~$\frac{2}{5}$\quad B.~$\frac{3}{5}$\quad C.~$\frac{5}{2}$\quad D.~$\frac{1}{5}$
\item Pour $f(x)=x^2+6$, quelle est la valeur de $f(-2)$ ?\par A.~$8$\quad B.~$24$\quad C.~$10$\quad D.~$2$
\end{enumerate}
\mxnewpage\subsection*{Partie 2 -- Exercices indépendants (14 points)}
Toute démarche doit être justifiée. Les résultats exacts sont attendus sauf indication contraire.
\mxneed{12}
\subsection*{Exercice 1 -- Une croissance par motif (5 points)}
Une famille de motifs est indexée par les entiers $n\geq1$. Le nombre de points du motif de rang $n$ est donné par $u_n=n^2+2n$. On cherche notamment le premier motif comprenant au moins 99 points ; aucune figure n’est nécessaire pour utiliser la formule donnée.
\begin{enumerate}
\item Calculer $u_1,u_2,u_3$. \hfill(1 pt)
\item Calculer $u_{n+1}-u_n$, puis établir le sens de variation de la suite $(u_n)$. \hfill(2 pts)
\item À partir de quel rang le motif contient-il au moins 99 points ? Justifier sans balayage infini. \hfill(2 pts)
\end{enumerate}
\mxneed{12}
\subsection*{Exercice 2 -- Deux dimensions et une contrainte (5 points)}
Un rectangle a un périmètre de 26 m et une aire de 40 mètres carrés. On note $x$ et $y$ ses deux dimensions, exprimées en mètres. Ces deux nombres sont strictement positifs. On comparera ensuite ce rectangle aux autres rectangles ayant le même périmètre.
\begin{enumerate}
\item Écrire les relations entre les deux longueurs x et y. \hfill(1 pt)
\item Déterminer les dimensions par une équation du second degré. \hfill(2 pts)
\item Parmi les rectangles de même périmètre, déterminer l’aire maximale et comparer à 40. \hfill(2 pts)
\end{enumerate}
\mxneed{12}
\subsection*{Exercice 3 -- Angle puis approximation (4 points)}
Les questions 1 et 2 concernent deux vecteurs $\vec u$ et $\vec v$ de normes respectives 2 et 3, tels que $\vec u\cdot\vec v=3$. La question 3 est indépendante : elle porte sur l’approximation par une tangente à la fonction carré.
\begin{enumerate}
\item Déterminer la mesure en degrés de l’angle non orienté entre $\vec u$ et $\vec v$. \hfill(1 pt)
\item Calculer la valeur exacte de $\norm{\vec u-\vec v}$. \hfill(1 pt)
\item Indépendamment, la tangente à la fonction carré en 1 a pour équation $y=2x-1$. L’utiliser pour approcher $\num{1,03}^2$, puis calculer la différence entre valeur exacte et approximation. \hfill(2 pts)
\end{enumerate}
